% Equation and Figure 4.2; requires \usepackage{graphicx}.
% If the thesis has its own subsection heading, do not repeat it here.
The relative effect of renormalisation-group evolution on an observable
$O$ is defined by
\begin{equation}
\delta_{\mathrm{RG}}(O)=100\,
\frac{O_{\mathrm{running}}-O_{\mathrm{frozen}}}
     {O_{\mathrm{frozen}}},
\label{eq:rg_correction_definition}
\end{equation}
where both calculations retain radiative threshold matching.
The frozen-parameter calculation disables parameter evolution, and
is not a tree-level neutrino-mass calculation.

\begin{figure}[htbp]
    \centering
    \includegraphics[width=0.95\textwidth]{Reports/thesis_chapter4/figures/fig_4_2_rg_corrections.pdf}
    \caption{Relative corrections from renormalisation-group evolution
    to (top) $\Delta m_{21}^{2}$ and (bottom) $\Delta m_{31}^{2}$
    for sixteen common-input T3 configurations, as defined in
    Eq.~\eqref{eq:rg_correction_definition}.
    Both running and frozen controls retain radiative threshold matching.
    A--E model families use consistent colours; the zero line
    separates enhancements and suppressions. Normal ordering is assumed,
    so the stored $\Delta m_{3\ell}^{2}$ is $\Delta m_{31}^{2}$.}
    \label{fig:rg_baseline_corrections}
\end{figure}
